\subsection{The reconstruction, its stage, and the five link diagrams}

\begin{result}{D.16}{Proved}{Canonical reconstruction from the supplied face and collar diagram}
Let $\mathcal V$ be the real oriented blow-up of a smooth complex normal-crossings Deligne--Mumford pair $(\mathcal X,\mathcal D)$, with its full deployed face diagram, collars, and maps to the strata. Collapsing precisely the normal angular circles over the zero radial coordinates recovers $\mathcal X$ as a topological stack. Equivalently, the mapping-cylinder pushout
\[
 \widehat X=X^*\mathop{\amalg}_{\dot X} C(p)
\]
of cote~150, scans~32--35, performed successively on the deployed strata, gives this same stack. For the supplied data, the final isomorphism is canonical and independent of the order of a compatible induction.

For $\overline{\mathcal M}_{g,n}$, the fundamental groupoid is already that of the compactification after all open divisors have been attached to the open stratum. In the deepest-strata-first reconstruction starting from $\mathcal V$, it is already that groupoid after the first layer, consisting of all maximal-corner fibres. In both descriptions it is the connected simply connected groupoid. These are two explicitly related filtrations of the same final quotient.
\end{result}
\begin{proof}
On a uniformizing chart, write the boundary equation as $z_1\cdots z_q=0$. The blow-up and its blow-down are
\[
 (\rho_i,\theta_i,z_{q+1},\ldots,z_m)\longmapsto
 (\rho_i e^{\mathrm i\theta_i},z_{q+1},\ldots,z_m),\qquad \rho_i\geq0.
\]
Its fibres identify exactly the $\theta_i$ for which $\rho_i=0$. The resulting quotient chart is the original complex chart: polar coordinates give a continuous bijection from the compact local quotient to its Hausdorff image, and the same assertion on relatively compact subcharts gives a homeomorphism. The construction is equivariant for the finite chart group and for coordinate changes. Taking the quotient groupoids therefore gives the original topological stack.

For a closed attaching stratum $S$, let $p:L_S\to S$ be its deployed normal link. Its mapping cylinder is
\[
 C(p)=\bigl(L_S\times[0,1]\amalg S\bigr)/((\ell,0)\sim p(\ell)).
\]
The collar identifies $L_S\times\{1\}$ with the attaching boundary of the previously cut piece. Gluing this cylinder implements exactly the preceding angular quotient on the corresponding normal cone. Repeating on all faces recovers the common quotient of $\mathcal V$. The coordinate equivalence relations commute: collapsing coordinate $i$ and then $j$ identifies exactly the same pairs as collapsing $j$ and then $i$. More generally the universal property of the quotient by the union of these relations identifies every compatible iterated pushout with one fixed quotient. The induced isomorphisms are unique over that quotient and satisfy their cocycle law. This is the asserted canonical independence for the given face/collar data; it requires no claim that collars have a preferred choice.

For the first filtration, set $\mathcal X^{\leq1}=\mathcal M_{g,n}\cup\{\text{open boundary divisors}\}$. Its inclusion into $\overline{\mathcal M}_{g,n}$ induces an isomorphism on $\pi_1$: the omitted strata have complex codimension at least two, and a loop or a homotopy of loops can be moved off them by real transversality in the uniformizing charts. Van Kampen for attaching an open divisor kills its normal meridian. Every essential simple closed curve gives a boundary degeneration, so the resulting quotient is $\Gamma_{g,n}$ modulo the normal subgroup generated by all Dehn twists. These twists generate the pure mapping class group. In positive genus this is exactly Labru\`ere--Paris, Corollary~2.11(i), with its puncture-induction proof on pages~86--89 \cite{D_LabruereParis2001}; in genus zero it follows from the pure-braid generators given in Part~B. Thus this quotient is trivial.

For the second filtration, the inclusion of the interior into $\mathcal V$ is a homotopy equivalence by pushing inward along a collar, so $\pi_1(\mathcal V)=\Gamma_{g,n}$. Over a maximal boundary stratum the deployed face is a torus bundle over the finite automorphism stack of a pants graph. The map to that stack is surjective on $\pi_1$, with kernel the image of its normal torus lattice. Replacing it by its mapping cylinder therefore imposes only the relations killing that lattice. Every essential curve extends to a pants decomposition; the normal closure of the lattices over all maximal strata is consequently the subgroup generated by all Dehn twists. The first complete layer has trivial fundamental group. Subsequent angular collapses again have connected fibres and induce quotients on $\pi_1$, so the group stays trivial. These intermediate quotients need not themselves be smooth stacks; they are the topological pushouts specified by the reconstruction. The two filtrations map to the same final quotient just proved.
\end{proof}

\paragraph{The four reconstruction cases.}
For $(0,4)$, $\mathcal V$ is a pair of pants. The three face maps are $S^1\to\{\text{point}\}$, and their mapping cylinders are disks. The pushout is $\mathbb P^1$. Killing the three peripheral loops already gives its trivial fundamental group.

For $(0,5)$, the ten deployed codimension-one faces are the normal circle bundles over the cut versions of the ten boundary copies of $\mathbb P^1$; their open bases are tripods. Fifteen deployed corner fibres are two-tori, one for each edge of the Petersen incidence graph. Each of the ten normal line bundles has degree $-1$. On a corner the two collapse maps are $(e^{\mathrm i\theta_1},e^{\mathrm i\theta_2})\mapsto e^{\mathrm i\theta_1}$ and $(e^{\mathrm i\theta_1},e^{\mathrm i\theta_2})\mapsto e^{\mathrm i\theta_2}$, followed by the corresponding remaining circle collapse. Their common quotient is the local chart $(\rho_1e^{\mathrm i\theta_1},\rho_2e^{\mathrm i\theta_2})$. Gluing these charts with the prescribed degree-$-1$ normal bundles recovers the degree-five del Pezzo surface $\overline M_{0,5}$. Collapsing all fifteen maximal two-torus fibres is already the first simply connected stage.

For $(1,1)$, the cut space is the compactified modular orbifold with its cusp cut open. The cusp face is $S^1\times B\mu_2\to B\mu_2$: the involution interchanges the two branches of the nodal rational curve and acts trivially on their tensor-product smoothing parameter. Its mapping cylinder fills the cusp while retaining that inertia. The resulting analytic stack is $\overline{\mathcal M}_{1,1}\simeq\mathbb P(4,6)$, including the interior inertia of orders $4$ and $6$ and generic inertia $\mu_2$. The boundary meridian maps to the usual Dehn twist $T\in\mathrm{SL}_2(\mathbb Z)$. The presentation $S^4=1$, $S^2=(ST)^3$ shows directly that setting $T=1$ kills the whole group. Thus the cusp attachment gives the compactification's trivial fundamental group; it retains the stack inertia.

For $(1,2)$, the two divisor normalizations are
\[
 [\overline M_{0,4}/\mathfrak S_2]\quad\hbox{and}\quad
 \overline{\mathcal M}_{1,1}\times\overline M_{0,3}.
\]
In the first, $\mathfrak S_2$ exchanges the node branches and fixes the two marked legs. The two maximal graph types are the loop with a rational tail and the two-edge cycle. Their deployed corner fibres are two-torus bundles over their graph-automorphism stacks; their local groups are respectively the groups $D$ and $C$ of B.10. The loop-tail corner has one retained normal direction for each divisor. The cycle corner has two distinct retained-branch incidences with the irreducible divisor, exchanged by its edge-swap automorphism. The face maps forget the smoothed angular coordinates and then use the sewing maps of B.11. Hence both incidences, and their finite automorphisms, occur in the pushout. The local two-coordinate quotient recovers the node-smoothing chart and the global quotient is $\overline{\mathcal M}_{1,2}$. Killing the two maximal meridian lattices kills every Dehn twist by the extension-to-pants argument; this is again the first simply connected stage.

\begin{result}{D.17}{Proved}{The five depth-three diagrams and the role of the four elementary types}
For $(0,5),(0,6),(1,2),(1,3),(2,0)$, the table below gives every boundary stratum of depth at most three, with its cut-component types and graph automorphisms. On the incidence category of deployed flags, the local group and each arrow are given by the following explicit normal-bundle and sewing construction. Its fundamental group is $\pi_1(L_{g,n})$. Consequently it equals $\Gamma_{g,n}$ in the three complexity-three types $(0,6),(1,3),(2,0)$; in the two complexity-two types its map to $\Gamma_{g,n}$ has kernel $F_{\aleph_0}$.

For every stable type, the action complex of groups on $P_2/\Gamma_{g,n}$ has fundamental group $\Gamma_{g,n}$. Its data use strata of dimension at most two. The connected one-dimensional stage supplies generators, the two-dimensional stage supplies relations, and a stratum of dimension three supplies relations among those relations. The four positive-dimensional single-component elementary types are exactly $(0,4),(1,1),(0,5),(1,2)$.
\end{result}
\begin{proof}
Here is the full geometric specification of the local groups, including their extension classes. For a graph $G$, let
\[
 \mathcal S_G=\left[\prod_v\mathcal M_{g_v,n_v}/\operatorname{Aut}(G)\right].
\]
For each edge $e=\{h,h'\}$ its normal line is
$\mathcal N_e=L_h^\vee\otimes L_{h'}^\vee$, of first Chern class
$-\psi_h-\psi_{h'}$. Let $\mathcal T_G$ be the simultaneous unit-circle bundle of these lines, with the graph group permuting the factors. Define $H_G=\pi_1(\mathcal T_G)$. Pulling to the product of cut-surface Teichm\"uller spaces makes this bundle trivial and proves
\[
 1\longrightarrow\mathbb Z^{E(G)}
 \longrightarrow H_G
 \longrightarrow\pi_1(\mathcal S_G)\longrightarrow1.
\]
The bundle itself specifies the extension, including its graph-automorphism lifts. Equivalently $H_G$ is the stabilizer of the chosen multicurve under sewing; opposite boundary twists of a cut edge are identified with the same global node twist.

A flag colours the edges by the stage at which they are contracted. Its group is the subgroup of $H_G$ preserving that colouring. The two edge injections into adjacent vertex groups are induced by sewing the specified boundary pairs and by retaining the other parametrized normal circles. On the node lattices they send a retained edge basis vector to that same edge basis vector. Graph automorphisms are restricted to those preserving the flag, before sewing. Thus multiple incidences are retained even when their contracted graph types coincide. All graph lists, individual contractions, and edge multiplicities are in the supplied JSON computation; the displayed table groups only isomorphic label-permutation patterns.

For completeness, the twisting elements in an ordinary presentation are specified as follows. Take the barycentric face category of the action complex, choose a representative of each orbit of a chain, and for an arrow $\alpha:A\to B$ choose $g_\alpha\in\Gamma_{g,n}$ carrying the chosen representative of $B$ to its designated face of the chosen representative of $A$. Then
\[
 \psi_\alpha(x)=g_\alpha^{-1}xg_\alpha,\qquad
 f_{\beta,\alpha}=g_\beta^{-1}g_\alpha^{-1}g_{\beta\alpha}
 \in H_C\quad(A\xrightarrow{\alpha}B\xrightarrow{\beta}C).
\]
The equality $\psi_\beta\psi_\alpha=\operatorname{Ad}(f_{\beta,\alpha})\psi_{\beta\alpha}$ follows by multiplying these three words. Multiplying four transports gives the cocycle identity, so this specifies the twisting on every triangle. The presentation has all local groups and arrow letters $u_\alpha$, the relations
\[
 u_\alpha x u_\alpha^{-1}=\psi_\alpha(x),\qquad
 u_\beta u_\alpha=f_{\beta,\alpha}u_{\beta\alpha},
\]
and $u_\alpha=1$ on a chosen maximal tree. This construction is the Borel action groupoid of the actual chain complex: equivariant simplexwise contraction gives its classifying-space comparison, so its homotopy fibre over $B\Gamma_{g,n}$ is that chain complex. For the depth-three link diagram D.2 supplies the one-type equivalence with $L_{g,n}$; D.1 proves exactly the group conclusions stated above.

For $P_2$, D.5 gives simple connectedness, and the same Borel fibre sequence gives $\pi_1=\Gamma_{g,n}$. If $d\leq2$, this assertion includes the empty multicurve, whose vertex is the open stratum. The exact map from the nonempty boundary diagram is induced by adjoining that empty multicurve and all incident chains; its kernel in the two complexity-two examples is the free group in D.1. Thus the low-dimensional and punctured-boundary presentations are related by this concrete map.

At a stratum of dimension one, the positive-complexity cut components have type $(0,4)$ or $(1,1)$, with pants factors. At dimension two they are either a single $(0,5)$ or $(1,2)$ component, or two complexity-one components. This follows by summing $3g_v-3+n_v$ and imposing stability, and proves that precisely the four elementary types furnish the local generator and relation data. D.5's connectedness of $P_1$ and simple connectedness of $P_2$ prove their global roles. At residual dimension three, the upper link is a join of curve complexes with total residual complexity three; its homotopy type is a wedge of two-spheres. Coning that link attaches three-dimensional cells up to homotopy and specifies identities among the two-dimensional relations. The single positive-complexity types are $(0,6),(1,3)$ and, only for the entire closed surface, $(2,0)$; the other possibilities are one complexity-two and one complexity-one component, or three complexity-one components. This lists every stratum type. For the three requested dimension-three moduli spaces, the sole dimension-three stratum is the open stratum itself. The two requested dimension-two spaces have no dimension-three stratum.
\end{proof}

The following enumeration was executed. Vertices are numbered from zero; $00$ denotes a loop and repeated entries denote parallel edges. The labelled number counts the individual strata in that label-permutation orbit. The normal-bundle rule, with these exact edges and cut types, specifies each vertex group and every coloured-flag subgroup.

{\small\paragraph{Type $(0,5)$.}
\begin{longtable}{c|r|l|l|l|r}
$q$ & labelled number & genera; leg counts & edge multiset & cut types & $|\operatorname{Aut}G|$\\\hline
1 & 10 & $(0,0);(2,3)$ & $\{01\}$ & $(0,3),(0,4)$ & 1\\
2 & 15 & $(0,0,0);(1,2,2)$ & $\{01,02\}$ & $(0,3),(0,3),(0,3)$ & 1\\
\end{longtable}
\paragraph{Type $(0,6)$.}
\begin{longtable}{c|r|l|l|l|r}
$q$ & labelled number & genera; leg counts & edge multiset & cut types & $|\operatorname{Aut}G|$\\\hline
1 & 15 & $(0,0);(2,4)$ & $\{01\}$ & $(0,3),(0,5)$ & 1\\
1 & 10 & $(0,0);(3,3)$ & $\{01\}$ & $(0,4),(0,4)$ & 1\\
2 & 60 & $(0,0,0);(1,2,3)$ & $\{01,02\}$ & $(0,3),(0,3),(0,4)$ & 1\\
2 & 45 & $(0,0,0);(2,2,2)$ & $\{01,02\}$ & $(0,3),(0,3),(0,4)$ & 1\\
3 & 15 & $(0,0,0,0);(0,2,2,2)$ & $\{01,02,03\}$ & $(0,3),(0,3),(0,3),(0,3)$ & 1\\
3 & 90 & $(0,0,0,0);(1,1,2,2)$ & $\{01,02,13\}$ & $(0,3),(0,3),(0,3),(0,3)$ & 1\\
\end{longtable}
\paragraph{Type $(1,2)$.}
\begin{longtable}{c|r|l|l|l|r}
$q$ & labelled number & genera; leg counts & edge multiset & cut types & $|\operatorname{Aut}G|$\\\hline
1 & 1 & $(0);(2)$ & $\{00\}$ & $(0,4)$ & 2\\
1 & 1 & $(0,1);(2,0)$ & $\{01\}$ & $(0,3),(1,1)$ & 1\\
2 & 1 & $(0,0);(0,2)$ & $\{00,01\}$ & $(0,3),(0,3)$ & 2\\
2 & 1 & $(0,0);(1,1)$ & $\{01,01\}$ & $(0,3),(0,3)$ & 2\\
\end{longtable}
\paragraph{Type $(1,3)$.}
\begin{longtable}{c|r|l|l|l|r}
$q$ & labelled number & genera; leg counts & edge multiset & cut types & $|\operatorname{Aut}G|$\\\hline
1 & 1 & $(0);(3)$ & $\{00\}$ & $(0,5)$ & 2\\
1 & 3 & $(0,1);(2,1)$ & $\{01\}$ & $(0,3),(1,2)$ & 1\\
1 & 1 & $(0,1);(3,0)$ & $\{01\}$ & $(0,4),(1,1)$ & 1\\
2 & 1 & $(0,0);(0,3)$ & $\{00,01\}$ & $(0,3),(0,4)$ & 2\\
2 & 3 & $(0,0);(1,2)$ & $\{00,01\}$ & $(0,3),(0,4)$ & 2\\
2 & 3 & $(0,0);(1,2)$ & $\{01,01\}$ & $(0,3),(0,4)$ & 2\\
2 & 3 & $(0,0,1);(1,2,0)$ & $\{01,02\}$ & $(0,3),(0,3),(1,1)$ & 1\\
3 & 3 & $(0,0,0);(0,1,2)$ & $\{00,01,12\}$ & $(0,3),(0,3),(0,3)$ & 2\\
3 & 3 & $(0,0,0);(0,1,2)$ & $\{01,01,02\}$ & $(0,3),(0,3),(0,3)$ & 2\\
3 & 1 & $(0,0,0);(1,1,1)$ & $\{01,02,12\}$ & $(0,3),(0,3),(0,3)$ & 1\\
\end{longtable}
\paragraph{Type $(2,0)$.}
\begin{longtable}{c|r|l|l|l|r}
$q$ & labelled number & genera; leg counts & edge multiset & cut types & $|\operatorname{Aut}G|$\\\hline
1 & 1 & $(1);(0)$ & $\{00\}$ & $(1,2)$ & 2\\
1 & 1 & $(1,1);(0,0)$ & $\{01\}$ & $(1,1),(1,1)$ & 2\\
2 & 1 & $(0);(0)$ & $\{00,00\}$ & $(0,4)$ & 8\\
2 & 1 & $(0,1);(0,0)$ & $\{00,01\}$ & $(0,3),(1,1)$ & 2\\
3 & 1 & $(0,0);(0,0)$ & $\{00,01,11\}$ & $(0,3),(0,3)$ & 8\\
3 & 1 & $(0,0);(0,0)$ & $\{01,01,01\}$ & $(0,3),(0,3)$ & 12\\
\end{longtable}
}

\paragraph{Concrete group forms in the first three examples.}
For $(0,5)$, each Petersen vertex $i$ has group
\[
 H_i=\langle f_i,h_{ij_1},h_{ij_2},h_{ij_3}\mid
 [f_i,h_{ij_a}]=1,\quad h_{ij_1}h_{ij_2}h_{ij_3}=f_i\rangle.
\]
This is $\mathbb Z\times F_2$; the relation records normal degree $-1$ and fixes the lifts of the three base peripheral loops. The edge group is $\mathbb Z^2$. Plumbing injects its basis as $(f_i,h_{ij})$ at $i$ and $(h_{ji},f_j)$ at $j$. Non-tree edges carry their arrow letters in the preceding presentation. Abelianizing gives the matrix $-I+A$ of D.4, and hence rank $11$, confirming that this link group is the specified extension of $\Gamma_{0,5}$.

For $(0,6)$, the fifteen $2|4$ divisor groups are the normal circle extensions of $\Gamma_{0,5}$, and the ten $3|3$ divisor groups those of $F_2\times F_2$, with Euler classes $-\psi_h-\psi_{h'}$. The $105$ two-edge groups are the two-circle extensions over a tripod, hence abstractly $\mathbb Z^2\times F_2$; their chosen peripheral lifts remain the prescribed normal-line lifts. The $105$ three-edge groups are $\mathbb Z^3$. All flag maps are the specified sewing maps, and the triangles are exactly the nested edge-contraction flags of the enumerated trees.

For $(1,2)$, B.10--B.11 give the completely presented node-only groups
\[
\begin{aligned}
 N&=\langle a,c,r\mid [a,c]=[a,r]=1,\ r^2=1\rangle,\\
 T&=\langle a,b\mid aba=bab\rangle,\\
 D&=\langle a,h\mid[a,h]=1\rangle,\\
 C&=\langle a,c,q\mid[a,c]=1,\ q^2=1,\ qaq^{-1}=c\rangle.
\end{aligned}
\]
Here $N,T$ are the irreducible and separating divisor groups; $D,C$ are the loop-tail and two-cycle groups. Their node lattices are $\langle a\rangle$, $\langle(ab)^6\rangle$, $\langle a,h^2\rangle$, and $\langle a,c\rangle$. The injections $D\to T$ and $D\to N$ are
\[
 a\mapsto a,\ h\mapsto(ab)^3;\qquad
 a\mapsto a,\ h\mapsto ac^{-1}r.
\]
For a branch of $C\to N$ the flag subgroup is $\langle a,c\rangle$ and its injection is the identity; the other branch is its $q$-conjugate. This restriction is necessary because $q$ exchanges the retained and smoothed edges. The map to the open presentation
$\Gamma_{1,2}=B_4(a,b,c)/\langle(abc)^4\rangle$
is
$h=(ab)^3$, $q=(abc)^2$, $r=(aba)^{-1}(abcaba)(aba)$.
The word checks and external-leg-framed versions are proved and executed in B.9--B.11.

The complete finite word dictionary for the genus-zero $P_2$ action presentations is proved in D.22, using B.18--B.19's exact braid normalization.  D.23 gives the entire fifteen-cell diagram for $(1,2)$.  The remaining named positive-genus transporters and dimension-three word identities are recorded explicitly as Open in D.24.  B.20--B.21 separately give the complete sixty-core-path presentation of the bipyramid--pyramid model.
